\chapter{总结与展望}\label{chp:summary}

\section{本文总结}

随着大语言模型在智能问答、知识检索、文本生成等场景中的广泛应用，模型生成内容的真实性、可靠性与可解释性问题日益突出。幻觉现象作为制约大语言模型安全落地的重要因素，已经成为当前生成式人工智能研究中的关键问题。因此，围绕大语言模型幻觉检测问题，探索能够充分利用模型内部表征、兼顾检测性能与方法可解释性的检测技术，具有重要的理论意义与应用价值。本文围绕“基于内部表征的大语言模型幻觉检测方法研究”这一主题，从通用幻觉检测与RAG幻觉检测两个方向展开研究，针对现有方法在内部表征利用不充分、外部知识信号建模不足等问题，提出了两种具有针对性的幻觉检测方法。

针对通用幻觉检测任务，本文提出了一种基于多方向表征建模的幻觉检测方法ATScore。该方法从大语言模型生成过程中的内部隐藏状态出发，构建统一的三维激活张量表示，并分别从layer方向与token方向对内部表征进行链式建模，在此基础上进一步提取链长度、转折角度等几何统计特征，构造无需额外训练的白盒幻觉检测信号。与现有多数仅沿单一方向、仅利用局部静态表示的内部表征方法相比，ATScore能够在更统一的视角下刻画大模型内部状态的动态演化结构，更充分地挖掘与幻觉相关的异常模式。实验结果表明，ATScore在多个公开数据集和多种评价指标上均取得了良好的检测效果，整体优于多种黑盒与白盒基线方法，表现出较强的稳定性与跨场景适用性。同时，消融实验进一步验证了多方向建模策略以及不同聚合方式对检测性能的影响，说明从内部表征演化轨迹中提取几何特征是一种有效且具有普适性的幻觉检测思路。

针对RAG幻觉检测任务，本文提出了一种基于检索知识冲突消解的幻觉检测方法SKID。该方法立足于RAG场景中“模型内部知识—外部检索知识”并存且可能发生冲突的基本特点，从大语言模型内部表征出发，首先借助稀疏自编码器将原始隐藏状态映射到高维稀疏特征空间，以获得更具可解释性的语义表征；随后通过互信息筛选获得与幻觉判别最相关的稀疏特征维度，并进一步结合检索文档对输出内容进行token级过滤，尽可能滤除受外部知识直接支持的生成成分，将检测重点集中于更能反映内部知识偏差的残差信号之上；最后使用可解释的广义加性模型对冲突信号进行建模，实现对RAG幻觉的有效检测。实验结果表明，SKID在多个RAG数据集和多种评价指标上均具有较好的竞争力，整体优于多类基线方法。进一步的分析还表明，SKID不仅能够提升RAG幻觉检测性能，而且能够通过稀疏特征的语义归纳揭示“模型为何被判定为产生幻觉”，从而在检测准确性与结果可解释性之间取得较好的统一。

总体来看，本文的两项研究工作在研究目标与技术路线方面具有清晰的递进关系。ATScore面向通用幻觉检测场景，重点解决现有方法对大模型内部表征利用不充分的问题，建立了从内部激活结构中挖掘幻觉信号的基本分析框架；SKID则进一步面向RAG幻觉检测场景，在继承内部表征分析思路的基础上，将外部检索知识显式纳入检测过程，通过冲突消解机制对内部知识与外部知识之间的关系进行建模，并进一步提升方法的可解释性。二者共同构成了一个从一般生成场景到检索增强生成场景、从内部动态建模到内外知识关系建模逐步推进的研究体系。

综上所述，本文围绕大语言模型幻觉检测问题，针对不同应用场景下的关键挑战，提出了两种基于内部表征的检测方法，并通过系统实验验证了其有效性、稳定性与一定的通用性。本文的研究工作表明，大语言模型内部表征中蕴含着丰富而可利用的幻觉判别信息；与此同时，在RAG场景下，若能进一步结合外部检索知识对内部表征进行冲突分析，则能够更准确地识别复杂生成条件下的幻觉行为。本文的研究不仅为后续大语言模型幻觉检测方法的设计提供了新的思路，也为构建更加安全、可信、可解释的大语言模型应用系统提供了有价值的方法支撑。

\section{未来展望}

尽管本文围绕通用幻觉检测与RAG幻觉检测两个方向进行了较为系统的研究，并取得了一定成果，但大语言模型幻觉问题本身仍具有高度复杂性与开放性，相关研究远未结束。随着模型规模、应用场景与外部工具使用方式的不断扩展，未来在幻觉检测领域仍有许多值得进一步深入探索的问题。结合本文已有工作，未来可从以下几个方面展开进一步研究。

首先，在通用幻觉检测方向上，仍可进一步提升对内部表征动态结构的建模能力。本文提出的ATScore方法主要从layer方向与token方向分别构造链式表征，并通过几何统计量刻画隐藏状态的演化特征。这一思路验证了内部表征动态变化中确实蕴含有价值的幻觉信号，但当前所使用的几何特征仍相对基础，主要集中于链长度、转折角度等较低阶统计信息。未来可以考虑从更丰富的链式表示出发，引入更加细粒度的动力学描述方式，例如局部波动模式、多尺度变化趋势、不同层级间的协同演化关系等，以进一步刻画模型在生成正确内容与生成幻觉内容时内部状态变化机制的差异。此外，当前方法主要在单模型、单次生成的设定下进行分析，未来还可以进一步探索其在跨模型迁移、跨任务泛化以及更复杂生成场景中的适用性，从而增强方法的普适性与鲁棒性。

其次，在RAG幻觉检测方向上，如何更精细地建模内部知识与外部检索知识之间的关系，仍然是一个值得持续深入的问题。本文提出的SKID方法通过过滤检索文档中已出现的token，并聚焦未被外部知识直接支持的生成片段，实现了对内外知识冲突的初步刻画。这种思路具有较好的直观性与有效性，但从更一般的角度来看，模型生成内容与检索知识之间的关系并不总是简单的“出现”或“未出现”关系。未来可以进一步考虑从事实蕴含、表征维度以及知识粒度对齐等更细致的层面，对外部知识支持关系进行建模。例如，某些生成内容虽然在表面token层面未直接出现在检索文档中，但其语义上可能仍受到检索证据支持；相反，也存在表面形式重合但事实关系发生偏移的情况。因此，如何由浅层的token匹配进一步走向更深层的语义冲突分析，是未来提升RAG幻觉检测准确性的重要方向。

综上所述，本文的研究工作为基于内部表征的大语言模型幻觉检测提供了有益探索，但无论是内部动态表征的深层建模、RAG场景下内外知识关系的细粒度刻画，还是检测结果的解释能力、在线反馈能力与系统化部署能力，仍然存在广阔的研究空间。未来，随着大语言模型能力的持续提升以及应用场景的不断拓展，面向高可靠性、高实用性与高可解释性的幻觉检测技术将具有更加重要的研究意义与应用价值。
